%% kapa: exam=1 round=21 year=2010 period=2 form=A subject=accounting subject_ko=회계학 seq=14 number=14
%% source: Q-Net 감정평가사 기출문제 (cst00309) artlSeq=5024181 file=extracted/1008/5024181/2010년 제21회 감정평가사 1차시험 기출문제/2교시.hwp
%% conversion: kordoc 4.24.51 -> tools/gen1.py (auto); review=proofread-gpt-sol; flags=none
\begin{question}{14}
20×0년 초 (주)감평은 조류발전을 목적으로 발전설비를 ￦240,000 현금으로 구입하여 설치완료하고 운영을 시작하였다. 이 발전설비는 내용연수 종료 후 원상복구의무가 있으며, 동 복구의무는 발전설비의 설치시점부터 충당부채의 인식요건을 충족시킨다. 내용연수 종료시점에 예상되는 복구비용은 ￦30,000으로 복구비용 산정 시 할인율은 연 10\%이다. 발전설비는 정액법(내용연수 4년, 잔존가치 ￦0)으로 감가상각한다. 동 발전설비 및 충당부채와 관련하여 (주)감평이 20×0년도에 인식할 감가상각비와 차입원가는 각각 얼마인가? (단, 10\%, 4기간의 ￦1에 대한 현가계수는 0.68301이며 할인율은 변동이 없고, 가장 근사치를 선택할 것)
\begin{choices}
  \item \makebox[16em][l]{감가상각비 ￦60,000}차입원가 ￦2,049
  \item \makebox[16em][l]{감가상각비 ￦60,000}차입원가 ￦2,798
  \item \makebox[16em][l]{감가상각비 ￦60,000}차입원가 ￦7,500
  \item \makebox[16em][l]{감가상각비 ￦65,123}차입원가 ￦2,049
  \item \makebox[16em][l]{감가상각비 ￦65,123}차입원가 ￦2,798
\end{choices}
\end{question}
\answer{4}
